% The Frohmanian Tether on the Coadjoint Orbit: A Lean-Synchronized Statement
% Copyright (c) 2026 Benjamin Stanley Frohman. Licensed under Apache-2.0.
% Engine: XeLaTeX (tectonic). Build: tectonic main.tex
\documentclass[11pt]{article}
\usepackage{frohman-paper}
\usepackage{fancyvrb}
\fvset{fontsize=\small,frame=single,framesep=4pt,rulecolor=\color{black!40}}
\renewcommand{\paperdoi}{10.5281/zenodo.23940666}
\renewcommand{\shorttitle}{The Frohmanian Tether on the Coadjoint Orbit}
\newcommand{\pin}{4aba751d2d365faf0cc48c5b46e9aa0113519273}
\newcommand{\pinshort}{4aba751}
\newcommand{\repo}{https://github.com/BenFrohman/frohmanian-tether-sdiff}
\newcommand{\src}[2]{\href{\repo/blob/\pin/#1\#L#2}{\file{#1:#2}}}
\newcommand{\tagbox}[2]{\textcolor{#1}{\textbf{\textsc{#2}}}}
\newcommand{\PROVED}{\tagbox{green!45!black}{proved}}
\newcommand{\COND}{\tagbox{orange!75!black}{conditional}}
\newcommand{\OPEN}{\tagbox{blue!60!black}{open}}
\newcommand{\REFUTED}{\tagbox{red!65!black}{refuted}}
\newcommand{\NF}{\tagbox{gray!80!black}{not formalized}}
\newtcolorbox{authortext}{enhanced,breakable,colback=yellow!4,colframe=black!45,boxrule=0.5pt,arc=1pt,left=6pt,right=6pt,top=4pt,bottom=4pt,title={\small Author's text, 4 November 2026 (verbatim)},fonttitle=\bfseries\small,coltitle=black,colbacktitle=yellow!18}
\newtcolorbox{leancorr}{enhanced,breakable,colback=red!3,colframe=red!55!black,boxrule=0.6pt,arc=1pt,left=6pt,right=6pt,top=4pt,bottom=4pt,fontupper=\small,title={\small [Lean correction --- the Lean development at \pinshort\ wins]},fonttitle=\bfseries\small}
\newtcolorbox{leanstatus}{enhanced,breakable,colback=green!2,colframe=green!35!black,boxrule=0.5pt,arc=1pt,left=6pt,right=6pt,top=4pt,bottom=4pt,fontupper=\small,title={\small Lean status at commit \pinshort},fonttitle=\bfseries\small}
\begin{document}
\thispagestyle{plain}
\begin{center}
{\small\scshape\color{red!60!black} Preprint --- not peer reviewed\par}\vspace{10pt}
{\LARGE\bfseries The Frohmanian Tether on the Coadjoint Orbit:\\[2pt] A Lean-Synchronized Statement\par}\vspace{10pt}
{\large\scshape Benjamin Stanley Frohman\par}\vspace{1pt}
{\small (B.\,S.~Frohman)\par}\vspace{3pt}
{\small Independent Researcher, Texas, United States\par}\vspace{2pt}
{\small ORCID \href{https://orcid.org/0009-0006-7068-3718}{0009-0006-7068-3718}\quad$\cdot$\quad\href{mailto:frohmanbenjamin@gmail.com}{\nolinkurl{frohmanbenjamin@gmail.com}}\par}\vspace{2pt}
{\small Version 1.0.0 $\cdot$ 11 November 2026 $\cdot$ DOI \doi{10.5281/zenodo.23940666}\par}
\end{center}
\vspace{4pt}

\begin{scopebox}
\textbf{Scope.} This is \emph{not} a Clay Millennium claim. Nothing here proves existence,
smoothness or global regularity of solutions of the 3D Navier--Stokes equations. It is a
statement of definitions and of what the Lean~4 development
\href{\repo}{\nolinkurl{BenFrohman/frohmanian-tether-sdiff}} proves at commit
\href{\repo/tree/\pin}{\texttt{\pinshort}}, item by item, with each item labelled
\PROVED, \COND\ (proved from explicitly named hypotheses), \OPEN, \REFUTED\ (its negation is
kernel-checked) or \NF. All Lean results depend only on the standard axioms
\texttt{propext}, \texttt{Classical.choice}, \texttt{Quot.sound}.
\end{scopebox}

\begin{abstract}
\noindent
On 4 November 2026 the author wrote a short draft, ``The Frohmanian Tether on the Coadjoint
Orbit'', synchronized by name with the Lean categories of the tether development:
\lean{CoadjointOrbit}, \lean{SDiffActionHyp}, \lean{TetherKernel}, \lean{TetheredBracket},
\lean{FunctionalDerivative}. This note reproduces that draft verbatim, section by section, and
places after each section the status of every claim in the current Lean development (commit
\texttt{\pinshort}, Lean \texttt{v4.30.0-rc1}, Mathlib \texttt{0c154d67}), with the exact Lean name
and \file{file:line}. Where the draft and the Lean disagree, the Lean wins, and the disagreement
is shown in a marked ``Lean correction'' box rather than silently rewritten. The main
corrections are: the field \lean{transitive_on_casimir_levels} was refuted and removed; the
identification $\delta H/\delta\omega = u$ is false for the formal derivative; fixed-functional
coadjoint invariance of the kernel is refuted for the concrete action (only a transported,
equivariant form is proved); $\Pi_u$ is an $L^2$ Gram--Schmidt projector, not a pointwise
one; the degeneracy condition C2 for \lean{TetherKernel} is not proved; and the tether rate, as
written, is refuted. The kernel definition with the global enstrophy weight, its symmetry,
its C3 quadratic form, its nondegeneracy and the restricted canonicity theorem are proved.
Periodic Beale--Kato--Majda continuation is now proved, so the main periodic theorem needs only the
open closing bound. A repaired rate \lean{FrohmanianTetherRateSmooth} yields a conditional reduction
of Fefferman's Problem~(A); both of its hypotheses remain open for general data.
\end{abstract}

\section*{How to read this note}
Each section opens with the author's text in a yellow box, unedited (inline code is set as
\lean{Lean} identifiers and the displayed Lean is reproduced exactly as written). The green box
that follows gives the Lean status. Red boxes mark places where the draft disagrees with the
Lean at \texttt{\pinshort}; the Lean statement is authoritative. File references link to
\href{\repo/tree/\pin}{the pinned commit}. Throughout, \lean{T3} is
\lean{EuclideanSpace ℝ (Fin 3)}, i.e.\ $\R^3$ with Lebesgue measure; it is \emph{not} the torus
(see \cref{rem:R3}). At this pin, periodic Beale--Kato--Majda continuation is \PROVED\
(\lean{periodicBKMContinuation_of_classical}, \file{Pillars/BKMContinuation.lean}); the only remaining
open hypothesis of the main periodic theorem is \lean{PeriodicTetherClosingBound} (as hard as Clay~(B)).

%=====================================================================
\section{Categories and basic objects}
\begin{authortext}
\begin{itemize}
\item \lean{T3} = \lean{EuclideanSpace ℝ (Fin 3)} (model space).
\item \lean{CoadjointOrbit} : structure with fields\\
  \lean{val : T3 → EuclideanSpace ℝ (Fin 3)}\\
  \lean{property : ∀ x, NavierStokes3D.div val x = 0}.
\item \lean{SDiff} : an abstract group equipped with a \lean{MulAction SDiff CoadjointOrbit}.
\item \lean{SDiffActionHyp SDiff} : hypothesis structure whose fields are\\
  \lean{enstrophy_invariant}, \lean{casimir_invariant}, \lean{transitive_on_casimir_levels}.\\
  (No global axioms; every theorem that needs them takes the structure as an argument.)
\item \lean{CoadjointAction g ω := g • ω} (the \lean{MulAction}).
\item \lean{FunctionalDerivative} : Gâteaux representative on the orbit (or 0 if none exists).
\item \lean{velocity_from_vorticity} : Biot–Savart recovery of the velocity (Fourier-side
  inversion proved mode-by-mode; real-space kernel still open).
\end{itemize}
\end{authortext}

\begin{leanstatus}
\begin{itemize}
\item \lean{NavierStokes3D.T3} is \lean{EuclideanSpace ℝ (Fin 3)}: \src{FrohmanianTether/NS_Equations.lean}{64}. \textbf{Matches.}
\item \lean{ArnoldGeometric.CoadjointOrbit} is a structure with exactly the fields \lean{val} and \lean{property} as written: \src{FrohmanianTether/ArnoldGeometric.lean}{167}. \textbf{Matches.}
\item The abstract group with \lean{[Group SDiff] [MulAction SDiff CoadjointOrbit]} and \lean{CoadjointAction g ω := g • ω} (an \lean{abbrev}): \src{FrohmanianTether/ArnoldGeometric.lean}{177}. Naturality \lean{CoadjointAction_one}, \lean{CoadjointAction_mul} \PROVED\ (lines 185, 189). \textbf{Matches.}
\item No \lean{axiom} declarations anywhere; the strict audit passes (\cref{sec:repro}). \textbf{Matches.}
\item \lean{ArnoldGeometric.FunctionalDerivative} (\src{FrohmanianTether/ArnoldGeometric.lean}{367}): the \emph{unique} functional gradient (\lean{isFunctionalGradient_unique}, line 353, \PROVED) when one exists, else $0$. A functional gradient (\lean{IsFunctionalGradient}, line 337) is an element of the divergence-free $L^2$ closure \lean{DivFreeL2} at a regular vorticity whose pairing with every compactly supported smooth test direction $\varepsilon$ equals $\tfrac{d}{dt}\big|_0 F(\omega+t\varepsilon)$. \textbf{Matches in substance}; see the note below.
\item \lean{ArnoldGeometric.velocity_from_vorticity} (line 1017) is the real-space Biot--Savart integral on $\R^3$ (\lean{velocity_from_vorticity_eq_BiotSavart}, line 1022). Fourier side: the symbol identities \lean{PlanSteps.bsSymbol_div} and \lean{PlanSteps.bsSymbol_curl} (\src{FrohmanianTether/PlanSteps.lean}{41}, 45) are \PROVED\ for each wavevector $k\neq 0$. Real-space $\curl\circ\BS=\mathrm{id}$ is \OPEN; $\Div \BS\omega = 0$ is \COND\ on an interchange hypothesis (\lean{FrohmanianTether.div_biot_savart_velocity_of_interchange}, \src{FrohmanianTether/SymplecticTether.lean}{492}).
\end{itemize}
\end{leanstatus}

\begin{leancorr}
\textbf{(1a) \lean{SDiffActionHyp} has one field, not three.} At \texttt{\pinshort} the structure
(\src{FrohmanianTether/ArnoldGeometric.lean}{219}) has the single field \lean{casimir_invariant}.
\lean{enstrophy_invariant} is a \emph{theorem} derived from it
(\lean{SDiffActionHyp.enstrophy_invariant}, line 225, \PROVED). The field
\lean{transitive_on_casimir_levels} was \emph{removed} because it is false for the concrete
measure-preserving action: \lean{ConcreteSDiff.not_transitive_on_casimir_levels}
(\src{FrohmanianTether/ConcreteSDiff.lean}{112}, \REFUTED; two constant fields have equal
Casimirs but no group element moves one to the other). The remaining hypothesis is satisfiable:
\lean{ConcreteSDiff.sdiffActionHyp} (line 98, \PROVED).

\textbf{(1b) ``Fourier-side inversion proved mode-by-mode''} is accurate only at the level of the
algebraic symbol, for each fixed $k\neq0$ with $k\cdot w=0$. No Fourier transform or Fourier
series inversion theorem is formalized.

\textbf{(1c) Gâteaux representative.} Gâteaux representatives are not unique in general
(\lean{KineticEnergyDerivative.not_unique_isGateauxRepresentative},
\src{FrohmanianTether/KineticEnergyDerivative.lean}{303}); \lean{FunctionalDerivative} is the
unique one in the divergence-free $L^2$ closure, which is what makes it well defined.
\end{leancorr}

%=====================================================================
\section{The tether kernel}
\begin{authortext}
\begin{Verbatim}
TetherKernel (ω : CoadjointOrbit) (F G : Functional) : ℝ :=
  -κ * (∫ x, ‖ω.val x‖ ^ 2) *
    ∫ x, ⟪Π_u (FunctionalDerivative F ω) x,
            Π_u (FunctionalDerivative G ω) x⟫
\end{Verbatim}
\begin{itemize}
\item Sign: \(-\kappa\) with \(\kappa>0\).
\item Weight: the global scalar \(\|\omega\|_{L^2}^2\) (overall enstrophy), not a pointwise density.
\item \(\Pi_u\) = pointwise projection orthogonal to \lean{velocity_from_vorticity ω}.
\end{itemize}
\end{authortext}

\begin{leanstatus}
\lean{FrohmanianTether.TetherKernel} (\src{FrohmanianTether/SymplecticTether.lean}{198}):
\begin{Verbatim}
TetherKernel (ω : CoadjointOrbit) (F G : Functional) : ℝ :=
  -kappa * Enstrophy ω * ∫ x,
      inner ℝ
        (Pi_u (velocity_from_vorticity ω) (FunctionalDerivative F ω) x)
        (Pi_u (velocity_from_vorticity ω) (FunctionalDerivative G ω) x)
    ∂volume
\end{Verbatim}
\begin{itemize}
\item Sign and weight \textbf{match}: \lean{Enstrophy ω} $=\int|\omega|^2$ is the global scalar
(\src{FrohmanianTether/ArnoldGeometric.lean}{207}; author decision of 3 November 2026 recorded in
the source). $\kappa$ is concrete: \lean{kappa_eq_three_halves} ($\kappa = C_{CZ}(3) = 3/2$) and
\lean{kappa_pos} (\src{FrohmanianTether/SymplecticTether.lean}{154}, 160), \PROVED.
\item \lean{Functional} is \lean{CoadjointOrbit → ℝ} (\src{FrohmanianTether/SymplecticTether.lean}{74}).
\end{itemize}
\end{leanstatus}

\begin{leancorr}
\textbf{(2a) $\Pi_u$ is not a pointwise projection.} \lean{ForMathlib.Pi_u}
(\src{FrohmanianTether/ForMathlib/Projection.lean}{35}) is the $L^2$ Gram--Schmidt projector onto
the complement of $\spn\{u\}$:
\[
  \Pi_u v \;=\; v - \frac{\int \langle v,u\rangle}{\int |u|^2}\,u ,
\]
with one global coefficient. It removes the $L^2$ component of $v$ along $u$; it does not make
$v(x)\perp u(x)$ at each point.
\end{leancorr}

\begin{remark}[\lean{T3} is $\R^3$]\label{rem:R3}
Because \lean{T3} is $\R^3$ with Lebesgue measure, a continuous, $\Z^3$-periodic, nonzero field
has non-integrable $|u|^2$ (\lean{ExplicitFlows.not_integrable_sq_of_periodic},
\src{FrohmanianTether/ExplicitFlows.lean}{664}, \PROVED). By Mathlib's convention a
non-integrable Bochner integral is $0$, so \lean{Enstrophy ω} $=0$ for every continuous periodic
nonzero vorticity, and hence \lean{TetherKernel ω F G} $=0$ there for all $F,G$. This is a
one-line consequence of the cited lemma and \lean{MeasureTheory.integral_undef}; it is not a
separately named Lean theorem. The kernel therefore carries no information about periodic
flows in the present formalization.
\end{remark}

%=====================================================================
\section{The tethered bracket}
\begin{authortext}
\begin{Verbatim}
TetheredBracket ω F G :=
  ClassicalLiePoissonBracket ω F G + TetherKernel ω F G
\end{Verbatim}
This is a \textbf{metriplectic} bracket:
\begin{itemize}
\item the classical summand is antisymmetric and satisfies Jacobi;
\item \lean{TetherKernel} is symmetric and negative-semidefinite;
\item the sum is neither antisymmetric nor Jacobi.
\end{itemize}
In particular \lean{TetheredBracket ω F F ≠ 0} in general.
\end{authortext}

\begin{leanstatus}
\lean{FrohmanianTether.TetheredBracket} (\src{FrohmanianTether/SymplecticTether.lean}{205}) is
\lean{ClassicalBracket F G ω + TetherKernel ω F G}.
\begin{itemize}
\item Classical summand antisymmetric: \lean{ArnoldGeometric.classicalBracket_antisymm}
(\src{FrohmanianTether/ArnoldGeometric.lean}{1211}) and \lean{classicalBracket_self} (line 1224), \PROVED.
\item \lean{TetherKernel} symmetric: \lean{FrohmanianTether.tetherKernel_comm}
(\src{FrohmanianTether/SymplecticTether.lean}{261}), \PROVED.
\item C3 as an equality, $\mathfrak T(F,F) = -\kappa\,\|\omega\|^2\int\|\Pi_u\delta F\|^2$:
\lean{tetherKernel_quadratic_form} and \lean{tetherKernel_C3} (lines 247, 292), \PROVED.
Negative semidefiniteness follows in one line from this, \lean{kappa_pos} and nonnegativity of
the integrals; there is no separately named lemma \lean{TetherKernel ω F F ≤ 0}.
\item Nondegeneracy: \lean{TetherNondegeneracy.tetherKernel_nondegenerate}
(\src{FrohmanianTether/TetherNondegeneracy.lean}{320}), \PROVED: there are a regular $\omega$ and a
test functional $F$ with \lean{TetherKernel ω F F ≠ 0}. With \lean{classicalBracket_self} this
gives \lean{TetheredBracket F F ω ≠ 0} for that pair, so the sum is not antisymmetric
(combination of two proved lemmas; not a separately named theorem).
\end{itemize}
\end{leanstatus}

\begin{leancorr}
\textbf{(3a) Name and argument order.} The Lean name is \lean{ClassicalBracket}
(\src{FrohmanianTether/ArnoldGeometric.lean}{1010}), not \lean{ClassicalLiePoissonBracket}, and
both brackets take $\omega$ last: \lean{ClassicalBracket F G ω}, \lean{TetheredBracket F G ω}.
(\lean{TetherKernel ω F G} takes $\omega$ first, as written.)

\textbf{(3b) Jacobi for the classical bracket is not proved.} Only the pointwise algebraic core is
proved: \lean{Symplectori.SymplecticConstraint.cross_jacobi_identity} and
\lean{kksPointwise_cyclic_jacobi_identity} (\src{FrohmanianTether/SymplecticConstraint.lean}{161},
183). The functional Jacobi identity for \lean{ClassicalBracket} is \OPEN. A former theorem
claiming Jacobi for the tethered bracket was circular and was deleted (source comment at
\src{FrohmanianTether/SymplecticTether.lean}{653}).

\textbf{(3c) ``Neither \dots\ Jacobi''} for the sum is \NF: no Lean counterexample to the Jacobi
identity of \lean{TetheredBracket} exists (\lean{jacobiator}, line 665, is only defined).

\textbf{(3d) ``Metriplectic''} is \COND. The usual metriplectic requirements include degeneracy of
the symmetric part on the Hamiltonian (C2), which is not proved for \lean{TetherKernel}
(see (4a)). The structural facts proved are: antisymmetric $+$ symmetric, with the symmetric part
satisfying C3.
\end{leancorr}

%=====================================================================
\section{Properties that hold (under the named hypotheses)}
\begin{authortext}
\begin{itemize}
\item Energy conservation: \lean{TetherKernel ω KineticEnergyHamiltonian _ = 0} (C2) implies the
tethered flow still conserves \(H\).
\item Route-2 canonicity: any bilinear form that saturates the quadratic bound and polarizes to a
symmetric \(Q\) depending on the derivatives only through \(\Pi_u\) coincides with
\lean{TetherKernel}. Locality is the one-line corollary of that representation. The proof uses
only saturation + polarization.
\item Coadjoint invariance of the kernel is conditional on \lean{SDiffActionHyp} (enstrophy
invariance + density equivariance).
\end{itemize}
\end{authortext}

\begin{leanstatus}
\begin{itemize}
\item C2 is the predicate \lean{DegenerateWRTKineticEnergy}
(\src{FrohmanianTether/SymplecticTether.lean}{230}): \lean{∀ F ω, B ω F KineticEnergyHamiltonian = 0}.
For \lean{TetherKernel} only the trivial case is proved:
\lean{tetherKernel_C2_of_no_gateaux} (line 635, \PROVED) --- if $H$ has no functional gradient
at $\omega$, the pairing vanishes. The classical part conserves $H$:
\lean{ArnoldGeometric.classicalBracket_H_H} (\src{FrohmanianTether/ArnoldGeometric.lean}{1230}),
\PROVED.
\item Route-2 canonicity:
\lean{FrohmanianTether.uniqueness_of_minimal_symmetric_tether}
(\src{FrohmanianTether/Uniqueness.lean}{273}), \COND\ and \PROVED\ as an implication: if $B$
is symmetric, satisfies \lean{SaturatesTetherQuadratic} (line 226) and
\lean{PolarizesOnDifferentiable} (line 241), then $B\,\omega\,F\,G = $ \lean{TetherKernel ω F G}
for all $F,G$ that have functional gradients at $\omega$. The proof uses only saturation and
polarization of $B$ and of the tether (\lean{tetherKernel_polarizesOnDifferentiable}, line 250,
\PROVED); its C1--C3 binders are carried but unused. \textbf{Matches} the draft's ``uses only
saturation + polarization''.
\item Factorization through $\Pi_u$: \lean{step2_quadratic} (line 411) proves
\lean{FactorsThroughProjection B} (one symmetric $Q$ per $\omega$) from symmetry and the named
hypothesis \lean{h_proj} ($B$ sees $\delta F$ only through $\Pi_u\delta F$, \OPEN\ for a general
$B$). Locality is the corollary \lean{step1_of_factorization} (line 445), packaged as
\lean{step2_quadratic_and_locality} (line 458), \PROVED. \textbf{Matches} ``locality is the
one-line corollary''.
\end{itemize}
\end{leanstatus}

\begin{leancorr}
\textbf{(4a) Energy conservation is conditional on C2, and C2 is open.} The former route to C2,
via \lean{FunctionalDerivative KineticEnergyHamiltonian ω = velocity_from_vorticity ω}, is false
with these definitions
(\lean{KineticEnergyDerivative.not_forall_functionalDerivative_kineticEnergy_eq_velocity},
\src{FrohmanianTether/KineticEnergyDerivative.lean}{286}, \REFUTED, by a rough indicator field),
and the lemmas that assumed it were deleted. No tethered flow (time evolution generated by
\lean{TetheredBracket}) is defined in Lean, so the conservation statement itself is \NF; it is a
one-line consequence of C2 and \lean{classicalBracket_H_H} once such a flow is defined.

\textbf{(4b) Canonicity holds only on differentiable functionals.} Unrestricted polarization of the
tether is false: \lean{FrohmanianTether.not_polarizes_tetherKernel}
(\src{FrohmanianTether/Uniqueness.lean}{258}, \REFUTED). The conclusion is stated for $F,G$ with
functional gradients at $\omega$, which are the only ones a derivative-based kernel can see.

\textbf{(4c) Coadjoint invariance is not ``conditional on \lean{SDiffActionHyp}''.} With fixed test
functionals, C1 for \lean{TetherKernel} is \emph{false} for the concrete action:
\lean{FrohmanianTether.not_tetherKernel_invariant_concrete}
(\src{FrohmanianTether/TetherC1.lean}{191}, \REFUTED), although \lean{SDiffActionHyp} holds for
that action. What is proved is the \emph{transported} (equivariant) form, in which $F,G$ are
pulled back along $g$, for the translation subgroup:
\lean{FrohmanianTether.tetherKernel_transported_invariant}
(\src{FrohmanianTether/TetherEquivariantC1.lean}{272}, \PROVED; it uses enstrophy invariance and
the transformation law of the functional derivative). For a general abstract action, C1 is the
named open Prop \lean{TetherKernelInvariant} (\src{FrohmanianTether/SymplecticTether.lean}{626},
\OPEN).
\end{leancorr}

%=====================================================================
\section{The C2/C3 obstruction (open, recorded)}
\begin{authortext}
Symmetry of \lean{TetherKernel} together with C2 forces
\begin{Verbatim}
TetherKernel ω M KineticEnergyHamiltonian = 0
\end{Verbatim}
for every functional \lean{M}. Hence
\begin{Verbatim}
deriv (fun t => M (ω t)) = ClassicalLiePoissonBracket ω M KineticEnergyHamiltonian
\end{Verbatim}
The tether contributes nothing. The formal cubic damping term is absent.\\
(The projection gap remains as a secondary obstruction: only an upper bound on
\(\int|\Pi_u\omega|^2\) is available.)
\end{authortext}

\begin{leanstatus}
\begin{itemize}
\item The first display is C2 itself (\lean{DegenerateWRTKineticEnergy}, read in the right slot);
by \lean{tetherKernel_comm} the left slot vanishes too. This is a one-line consequence of a proved
lemma and the C2 hypothesis; it is not a separately named theorem. As a statement it is
\COND\ on C2, which is \OPEN\ for \lean{TetherKernel} (see (4a)).
\item The vanishing of the damping contribution of energy-null symmetric corrections is the subject
of the related record ``A Tether That Holds Nothing'' (\doi{10.5281/zenodo.23971661}).
\end{itemize}
\end{leanstatus}

\begin{leancorr}
\textbf{(5a)} The time-derivative identity is \NF: there is no Lean definition of the tethered
evolution $t\mapsto\omega(t)$, and the classical bracket is \lean{ClassicalBracket M H ω}
($\omega$ last). \textbf{(5b)} The ``projection gap'' remark is \NF\ in the repository; it is
recorded here as the author's mathematical observation only.
\end{leancorr}

%=====================================================================
\section{Analytic status}
\begin{authortext}
The differential inequality that would feed the Riccati comparison, the tether rate, and the
tether closing bound are \textbf{named open assumptions}. Any conditional regularity statement is
explicitly conditional on a resolution of the C2/C3 incompatibility. No Navier–Stokes existence,
uniqueness or regularity theorem is claimed.
\end{authortext}

\begin{leanstatus}
\begin{itemize}
\item Periodic BKM continuation is \PROVED:
\lean{FrohmanianTether.Pillars.periodicBKMContinuation_of_classical}
(\src{FrohmanianTether/Pillars/BKMContinuation.lean}{144}): $0<\nu$ implies
\lean{PeriodicBKMContinuation ν}.
\item Main theorem with BKM discharged:
\lean{FrohmanianTether.Pillars.clay_B_exact_of_periodic_tether_closing_bound}
(\src{FrohmanianTether/Pillars/BKMContinuation.lean}{170}) and
\lean{clay_B_exact_unique_of_periodic_tether_closing_bound} (line 177), \COND\ on the single named
hypothesis \lean{PeriodicTetherClosingBound} (\src{FrohmanianTether/ClayStatements.lean}{575}, \OPEN;
a time-uniform $\|\omega\|_\infty$ ceiling, by BKM as hard as Clay~(B) itself). This is not a proof
of Clay~(B). The two-hypothesis form
\lean{ClayStatements.clay_B_exact_of_periodic_tether_bound} (line 599) remains, and the discharged
form supplies \lean{hBKM} from the pillar.
\item PR~\#4 (merge commit \texttt{0698d08}, included in this pin) added the repaired rate
\lean{FeffermanClayA.FrohmanianTetherRateSmooth}
(\src{FrohmanianTether/FeffermanClayA.lean}{81}, \OPEN\ for general data) and the conditional theorem
\lean{FeffermanClayA.clay_A_fefferman_of_frohmanianTetherRateSmooth} (line 121):
if \lean{FrohmanianTetherRateSmooth} $u_0\,\nu$ holds for every $\nu>0$ and every Fefferman
datum $u_0$, and \lean{R3BKMContinuation} $\nu$ holds for every $\nu>0$, then
\lean{Clay_A_Fefferman}.
Both hypotheses remain open for general data and are only shown satisfiable at $u_0=0$
(\lean{frohmanianTetherRateSmooth_zero}, line 176; \lean{r3BKMContinuationAt_zero}, line 146).
This is not a proof of Clay~(A). The unrestricted \lean{FrohmanianTetherRate} remains \REFUTED\
(see (6a)).
\item \lean{FrohmanianTether.ClayStatements.FrohmanianTetherClosingBound} (line 48) is \OPEN.
\item The scalar ODE comparison machinery (\lean{TetheredLyapunov.comparison_ode_exists} and
neighbours, \file{FrohmanianTether/TetheredLyapunov.lean}) is \PROVED; the differential
inequality that would feed it for Navier--Stokes is not.
\end{itemize}
\end{leanstatus}

\begin{leancorr}
\textbf{(6a) The tether rate, as written, is false.}
\lean{TetherRateCounterexample.not_forall_frohmanianTetherRate}
(\src{FrohmanianTether/TetherRateCounterexample.lean}{428}, \REFUTED): a steady finite-energy
field with a nowhere-differentiable Takagi component satisfies the formal PDE through Mathlib's
zero junk derivative while its vorticity is unbounded. On periodic data it holds only vacuously
(\lean{TetherRatePeriodicVacuity.frohmanianTetherRate_of_periodic},
\src{FrohmanianTether/TetherRatePeriodicVacuity.lean}{29}). It is therefore not a usable open
assumption in its written form.

\textbf{(6b) Uniqueness is now proved.} Since the draft, uniqueness of classical $\Z^3$-periodic
smooth solutions has been formalized and imported as a pillar:
\lean{FrohmanianTether.Pillars.periodicSmoothSolutionOn_unique}
(\src{FrohmanianTether/Pillars/Uniqueness.lean}{65}, \PROVED, from
\href{https://github.com/BenFrohman/periodic-navier-stokes-smooth-uniqueness}{\nolinkurl{periodic-navier-stokes-smooth-uniqueness}}
at \texttt{de92662}). This is a classical result. No existence or regularity theorem is claimed.

\textbf{(6c)} The main conditional theorem above does not take a C2/C3 hypothesis; its only
remaining hypothesis is the closing bound (BKM continuation is now \PROVED). ``Conditional on a
resolution of the C2/C3 incompatibility'' applies to any \emph{damping} reading of the tether, not
to that theorem.
\end{leancorr}

%=====================================================================
\section{Open assumptions (all named, none axiomatic)}
\begin{authortext}
\begin{itemize}
\item \lean{SDiffActionHyp.transitive_on_casimir_levels}
\item \lean{FunctionalDerivative KineticEnergyHamiltonian ω = velocity_from_vorticity ω}
\item Euler equations realized as coadjoint motion (\lean{IsCoadjointMotion})
\item the tether rate and the tether closing bound
\item resolution of the C2/C3 obstruction (if a damping reading is wanted)
\end{itemize}
\end{authortext}

\begin{leanstatus}
\begin{itemize}
\item \lean{PlanSteps.IsCoadjointMotion} (\src{FrohmanianTether/PlanSteps.lean}{61}) is a named
open Prop, \OPEN. Under it and \lean{SDiffActionHyp}, every Casimir is conserved:
\lean{PlanSteps.casimir_conserved} (line 67), \COND.
\item Tether closing bound: \lean{PeriodicTetherClosingBound}, \OPEN\ (the only remaining open
hypothesis of the main periodic theorem). Periodic BKM continuation is \PROVED\ (see the analytic
status box).
\item C2 for \lean{TetherKernel}: \OPEN. C1 for an abstract action: \lean{TetherKernelInvariant}, \OPEN.
\item None is an \lean{axiom}; the audit reports zero axiom declarations.
\end{itemize}
\end{leanstatus}

\begin{leancorr}
\textbf{(7a)} \lean{transitive_on_casimir_levels} is no longer an open assumption: it is
\REFUTED\ for the concrete action and removed (see (1a)).
\textbf{(7b)} \lean{FunctionalDerivative KineticEnergyHamiltonian ω = velocity_from_vorticity ω}
is \REFUTED\ as a universal statement (see (4a)); the smooth-field identity is not decided.
\textbf{(7c)} The tether rate is \REFUTED\ as written (see (6a)); only the closing bound remains open.
\end{leancorr}

\begin{authortext}
All terminology, variable names and categories above match the current Lean development
(\lean{CoadjointOrbit}, \lean{SDiffActionHyp}, \lean{MulAction}, \lean{TetherKernel},
\lean{TetheredBracket}, \lean{FunctionalDerivative}, etc.). The geometric layer builds; the
dissipative bridge does not.
\end{authortext}

\begin{leancorr}
\textbf{(8a)} Names match except \lean{ClassicalLiePoissonBracket} (Lean: \lean{ClassicalBracket})
and the fields of \lean{SDiffActionHyp} (see (1a), (3a)). \textbf{(8b)} At \texttt{\pinshort} the whole
library builds and passes the strict axiom audit in CI (\cref{sec:repro}); the dissipative
content and the remaining analytic hypothesis \lean{PeriodicTetherClosingBound} enter only as named
open Props. Periodic BKM continuation is now proved.
\end{leancorr}

%=====================================================================
\section{Status table}
\small
\begin{longtable}{@{}>{\raggedright\arraybackslash}p{0.30\textwidth}>{\raggedright\arraybackslash}p{0.42\textwidth}>{\raggedright\arraybackslash}p{0.20\textwidth}@{}}
\toprule
\textbf{Item} & \textbf{Lean name (\file{file:line})} & \textbf{Status}\\\midrule\endhead
Model space $\R^3$ & \lean{NavierStokes3D.T3} (NS\_Equations:64) & definition\\
Orbit structure & \lean{ArnoldGeometric.CoadjointOrbit} (ArnoldGeometric:167) & definition\\
Action naturality & \lean{CoadjointAction_one}, \lean{CoadjointAction_mul} (185, 189) & \PROVED\\
Enstrophy invariance from Casimirs & \lean{SDiffActionHyp.enstrophy_invariant} (225) & \COND\\
Concrete action satisfies hyp. & \lean{ConcreteSDiff.sdiffActionHyp} (ConcreteSDiff:98) & \PROVED\\
Transitivity on Casimir levels & \lean{ConcreteSDiff.not_transitive_on_casimir_levels} (112) & \REFUTED\\
Uniqueness of functional gradient & \lean{isFunctionalGradient_unique} (ArnoldGeometric:353) & \PROVED\\
BS symbol: div-free, curl inverse & \lean{PlanSteps.bsSymbol_div}, \lean{bsSymbol_curl} (41, 45) & \PROVED\\
Real-space $\curl\circ\BS=\mathrm{id}$ & --- & \OPEN\\
Kernel definition, global weight & \lean{FrohmanianTether.TetherKernel} (SymplecticTether:198) & definition\\
$\kappa = 3/2 > 0$ & \lean{kappa_eq_three_halves}, \lean{kappa_pos} (154, 160) & \PROVED\\
Kernel symmetric & \lean{tetherKernel_comm} (261) & \PROVED\\
C3 quadratic form (equality) & \lean{tetherKernel_quadratic_form}, \lean{tetherKernel_C3} (247, 292) & \PROVED\\
Kernel nondegenerate & \lean{TetherNondegeneracy.tetherKernel_nondegenerate} (TetherNondegeneracy:320) & \PROVED\\
Classical bracket antisymmetric & \lean{ArnoldGeometric.classicalBracket_antisymm} (1211) & \PROVED\\
$\{H,H\}_{\mathrm{cl}}=0$ & \lean{classicalBracket_H_H} (1230) & \PROVED\\
Classical functional Jacobi & (pointwise only: \lean{cross_jacobi_identity}, SymplecticConstraint:161) & \OPEN\\
Tethered bracket not Jacobi & \lean{jacobiator} (665), defined only & \NF\\
C2 for the kernel & \lean{DegenerateWRTKineticEnergy} (230); trivial case \lean{tetherKernel_C2_of_no_gateaux} (635) & \OPEN\\
$\delta H/\delta\omega = u$ for all $\omega$ & \lean{not_forall_functionalDerivative_kineticEnergy_eq_velocity} (KineticEnergyDerivative:286) & \REFUTED\\
Energy conservation of tethered flow & --- (flow not defined) & \NF; \COND\ on C2\\
Restricted canonicity & \lean{uniqueness_of_minimal_symmetric_tether} (Uniqueness:273) & \COND\\
Tether polarizes on differentiable $F,G$ & \lean{tetherKernel_polarizesOnDifferentiable} (250) & \PROVED\\
Unrestricted polarization & \lean{not_polarizes_tetherKernel} (258) & \REFUTED\\
Factorization through $\Pi_u$ & \lean{step2_quadratic} (411) & \COND\ (on \lean{h_proj})\\
Locality corollary & \lean{step1_of_factorization}, \lean{step2_quadratic_and_locality} (445, 458) & \PROVED\\
C1, fixed functionals, concrete action & \lean{not_tetherKernel_invariant_concrete} (TetherC1:191) & \REFUTED\\
C1, transported, translations & \lean{tetherKernel_transported_invariant} (TetherEquivariantC1:272) & \PROVED\\
C1, abstract action & \lean{TetherKernelInvariant} (SymplecticTether:626) & \OPEN\\
C2/C3 obstruction (first display) & one line from \lean{tetherKernel_comm} $+$ C2 & \COND\ on C2\\
Tethered time-derivative identity & --- & \NF\\
Euler as coadjoint motion & \lean{PlanSteps.IsCoadjointMotion} (PlanSteps:61) & \OPEN\\
Casimir conservation & \lean{PlanSteps.casimir_conserved} (67) & \COND\\
Tether rate as written & \lean{not_forall_frohmanianTetherRate} (TetherRateCounterexample:428) & \REFUTED\\
Tether closing bound & \lean{PeriodicTetherClosingBound} (ClayStatements:575) & \OPEN\\
Periodic BKM continuation & \lean{periodicBKMContinuation_of_classical} (Pillars/BKMContinuation:144) & \PROVED\\
Clay (B) from bound (BKM discharged) & \lean{clay_B_exact_of_periodic_tether_closing_bound} (Pillars/BKMContinuation:170) & \COND\ (only \lean{hclose})\\
Clay (B) unique, same hypothesis & \lean{clay_B_exact_unique_of_periodic_tether_closing_bound} (177) & \COND\ (only \lean{hclose})\\
Two-hypothesis form (BKM as hyp.) & \lean{clay_B_exact_of_periodic_tether_bound} (ClayStatements:599) & \COND\\
Repaired rate (smooth class) & \lean{FrohmanianTetherRateSmooth} (FeffermanClayA:81) & \OPEN\\
Clay (A) from repaired rate $+$ $\R^3$ BKM & \lean{clay_A_fefferman_of_frohmanianTetherRateSmooth} (FeffermanClayA:121) & \COND\ (both hyps.\ open; shown at $u_0=0$)\\
Periodic uniqueness (classical) & \lean{Pillars.periodicSmoothSolutionOn_unique} (Pillars/Uniqueness:65) & \PROVED\\
Periodic fields not in $L^2(\R^3)$ & \lean{ExplicitFlows.not_integrable_sq_of_periodic} (ExplicitFlows:664) & \PROVED\\
\bottomrule
\end{longtable}
\normalsize
All file paths are under \file{FrohmanianTether/} at commit \texttt{\pinshort}.

%=====================================================================
\section{Related works}
This note is a standalone statement of the author's 4 November draft. It is \emph{not} a
re-release of any of the following records, which it cross-references:
\begin{itemize}
\item B.\,S.~Frohman, \emph{The Frohmanian Tether Bound Implies Global Smoothness on $\T^3$ under
Beale--Kato--Majda Continuation, and Implies BKM Integrability: Kernel-Checked Conditional Theorems
in Lean~4}, concept DOI \doi{10.5281/zenodo.24081999} (latest published version v0.1.5,
\doi{10.5281/zenodo.24122596}). The Lean development pinned here is commit \texttt{\pinshort}, which
includes the BKM continuation pillar.
\item B.\,S.~Frohman, \emph{The Frohmanian Tether on $\mathrm{SDiff}(\T^3)$: A Metriplectic
Correction, Its Conditional Canonicity, and an Obstruction to Damping}, Zenodo, 2026,
\doi{10.5281/zenodo.23956935}.
\item B.\,S.~Frohman, \emph{The Frohmanian Symplectic Tether $\mathfrak T$: Written Proofs Companion
to the Lean~4 Formalization}, v1.1.0, Zenodo, 2026, \doi{10.5281/zenodo.23942463}.
\item B.\,S.~Frohman, \emph{A Tether That Holds Nothing: Energy-Null Metriplectic Corrections and
the Vanishing Enstrophy Contribution on $\T^3$}, Zenodo, 2026, \doi{10.5281/zenodo.23971661}.
\end{itemize}

%=====================================================================
\section{Reproducibility}\label{sec:repro}
\begin{itemize}
\item Repository: \href{\repo}{\nolinkurl{https://github.com/BenFrohman/frohmanian-tether-sdiff}},
commit\\ \href{\repo/tree/\pin}{\texttt{\pin}} (\texttt{main}, 11 November 2026).
A source archive of this commit (\file{frohmanian-tether-sdiff-4aba751.zip}) is attached to this record.
\item Toolchain: \texttt{leanprover/lean4:v4.30.0-rc1}.
\item Mathlib: \texttt{0c154d67103f74be3a0f2c509f72ccbf5be9f2a7} (from \file{lake-manifest.json}).
\item CI for this commit: workflow ``Lean certificate'', run
\href{https://github.com/BenFrohman/frohmanian-tether-sdiff/actions/runs/43107567565}{43107567565}
(success; ends with \texttt{VERIFY OK}), and workflow ``Lean 4'', run
\href{https://github.com/BenFrohman/frohmanian-tether-sdiff/actions/runs/43107567433}{43107567433}
(success). The pull-request certificate for the same tree is run
\href{https://github.com/BenFrohman/frohmanian-tether-sdiff/actions/runs/43105941627}{43105941627}
(identical axiom artifact).
\item Axiom output (CI artifact \texttt{axioms-push-43107567565}, attached as
\file{axioms.txt} and \file{print_axioms.txt}):
\begin{Verbatim}
AUDIT SUMMARY: theorems audited: 1886 (1674 stated in the sources,
  212 auto-generated equation/congruence/injectivity lemmas);
  declarations audited: 2546; violations: 0
AUDIT PASS: only propext, Classical.choice, Quot.sound;
  no sorryAx; no axiom declarations
\end{Verbatim}
Every Lean theorem cited in this note appears in \file{axioms.txt} with
\texttt{[propext, Classical.choice, Quot.sound]}.
\item Rebuild (one command, from a clone at the pinned commit):
\begin{Verbatim}
git clone https://github.com/BenFrohman/frohmanian-tether-sdiff \
  && cd frohmanian-tether-sdiff \
  && git checkout 4aba751d2d365faf0cc48c5b46e9aa0113519273 \
  && bash scripts/verify.sh
\end{Verbatim}
\file{scripts/verify.sh} fetches the Mathlib cache for the pinned revision, runs \texttt{lake build},
the no-\texttt{sorry} check, the strict axiom audit and \texttt{\#print axioms}, and prints
\texttt{VERIFY OK}.
\end{itemize}

\section*{Provenance}
The author's text in the yellow boxes was written on 4 November 2026 and is reproduced without
change; the raw copy is attached as \file{original_draft.md}. Status boxes, corrections and the
table were checked against the Lean sources at \texttt{\pinshort} on 11 November 2026.

\vspace{1em}
\noindent\small\copyright\ 2026 Benjamin Stanley Frohman. Licensed under the Apache License,
Version 2.0. This document reflects the Lean development at the stated commit; it is not peer reviewed.
\end{document}
